\section{Related work}

\paragraph{Multilingual and cultural evaluation.}
Multilingual benchmark work has expanded language coverage and highlighted
large performance gaps across languages. MMLU-ProX evaluates advanced reasoning
across typologically diverse languages \citep{xuan2025mmluprox}, BenchMAX
proposes a broad multilingual evaluation suite \citep{huang2025benchmax}, and
OMGEval evaluates open-ended multilingual generation with localized questions
\citep{liu2024omgeval}. Global MMLU shows that simply translating English
benchmarks can preserve Western-centric cultural assumptions and distort model
rankings \citep{singh2024globalmmlu}. Cultural benchmarks such as CDEval and
MCEval argue that model behavior should be evaluated through culturally specific
and cross-cultural lenses \citep{wang2023cdeval,huang2025mceval}. Our work is
smaller and narrower: it does not attempt broad cultural coverage, but instead
measures the interactional cost of recovery after violating multilingual
interaction contracts.

\paragraph{Instruction following and code-switching.}
IFEval introduced verifiable instruction-following checks for LLMs
\citep{zhou2023ifeval}; XIFBench and Marco-Bench-MIF extend this concern to
multilingual instruction following with localized constraints
\citep{li2025xifbench,zeng2025marcomif}. OLA is the closest prior work to ours:
it studies output-language alignment in code-switched interactions and shows
that models often infer the wrong response language \citep{oh2026ola}. CodeMixQA
and related work further show that code-switched reasoning and generation remain
challenging \citep{winata2026codemixqa}. Re-prompt Tax includes output-language
alignment, but adds repair trajectories and a broader contract containing
script, preservation, register, and locale constraints, and standardized
repairs are needed to compare those failures as trajectories rather than
isolated first-turn answers.

\paragraph{Multi-turn dialogue, real-world chat logs, and judges.}
MT-Bench and Chatbot Arena evaluate conversational quality and preference using
multi-turn prompts and LLM or human judgments \citep{zheng2023mtbench}. WildChat
and LMSYS-Chat-1M provide large-scale public records of real LLM interactions
\citep{zhao2024wildchat,zheng2023lmsyschat1m}. We use WildChat only for a
bounded aggregate cue scan and release synthetic stress items instead of raw
user conversations. Our judge audit is related to LLM-as-a-judge methodology,
but multilingual judging itself is known to be inconsistent across languages
\citep{fu2025multilingualjudge}. We therefore treat the judge audit as a sanity
check rather than as final human validation.

\paragraph{Cost and accessibility.}
Prior work on multilingual tokenization shows that language users can face
unequal token costs and utility when using commercial APIs
\citep{ahia2023languagecost,petrov2023tokenizers}. Re-prompt Tax is
complementary: instead of asking whether the same content costs more to encode
in different languages, we ask how many extra turns and tokens are consumed when
the model violates a multilingual interaction contract.
